
Pendahuluan Jilid 1 \vtmpb{30.9.02}\pagereference{7}{7}

\medskip

Bab 11:  Ruang Ukur

\chapintrosection{10.9.96}{9}{9}

\section{111}{Aljabar-$\sigma$}{26.1.05}{10}{10}
{Definisi aljabar-$\sigma$;  himpunan terhitung;
aljabar-$\sigma$ yang dibangkitkan oleh suatu keluarga himpunan;  
aljabar-$\sigma$ Borel.}

\section{112}{Ruang ukur}{20.2.05}{14}{11}
{Definisi ruang ukur;  penggunaan $\infty$;  sifat-sifat
elementer;  himpunan terabaikan;  ukuran yang tertumpu pada
titik;  ukuran citra.}

\section{113}{Ukuran luar dan konstruksi \Caratheodory}
{6.4.05}{19}{13}
{Ukuran luar;  konstruksi \Caratheodory untuk membentuk suatu ukuran dari
ukuran luar.}

\section{114}{Ukuran Lebesgue pada $\Bbb{R}$}{14.6.05}{23}{13}
{Selang setengah terbuka;  ukuran luar Lebesgue;  ukuran Lebesgue;
himpunan Borel terukur.}

\section{115}{Ukuran Lebesgue pada $\BbbR^r$}{21.7.05}{28}{14}
{Selang setengah terbuka;  ukuran luar Lebesgue;  ukuran Lebesgue;
himpunan Borel terukur.}

\medskip

 Bab 12:  Integrasi

\chapintrosection{7.4.05}{35}{16}

\section{121}{Fungsi terukur}{21.12.03}{35}{16}
{Aljabar-$\sigma$ subruang;  fungsi bernilai riil yang terukur;
fungsi yang terdefinisi secara parsial;  fungsi terukur Borel;  operasi pada
fungsi terukur;  membangkitkan himpunan Borel dari setengah-ruang.}

\section{122}{Definisi integral}{4.1.04}{43}{18}
{Fungsi sederhana;  fungsi nonnegatif yang terintegralkan;  fungsi
bernilai riil yang terintegralkan;  fungsi yang terukur pada suatu himpunan
koterabaikan;  linearitas
integral.}

\section{123}{Teorema-teorema konvergensi}{18.11.04}{52}{20}
{Teorema B.Levi;  lema Fatou;  Teorema Konvergensi
Terdominasi Lebesgue;  mendiferensialkan di bawah tanda integral.}

\medskip

 Bab 13:  Pelengkap

\chapintrosection{16.6.01}{56}{21}

\section{131}{Subruang terukur}{18.3.05}{56}{21}
{Ukuran subruang pada himpunan bagian terukur;  integrasi pada
himpunan bagian terukur.}

\section{132}{Ukuran luar dari ukuran}{6.4.05}{58}{22}
{Ukuran luar yang terkait dengan suatu ukuran;  ukuran luar Lebesgue
sekali lagi;  selubung terukur.}

\section{133}{Konsep integrasi yang lebih luas}{29.3.10}{61}{23}
{$\infty$ sebagai nilai suatu integral;  fungsi bernilai kompleks;
integral atas dan bawah.}

\section{134}{Lebih lanjut tentang ukuran Lebesgue}{7.1.04}{68}{25}
{Invariansi translasi;  himpunan tak terukur;  regularitas dalam dan luar;
himpunan dan fungsi Cantor;  *integral Riemann.}

\section{135}{Garis bilangan riil diperluas}{14.9.04}{79}{27}
{Aljabar $\pm\infty$;  himpunan Borel dan barisan konvergen
dalam $[-\infty,\infty]$;   fungsi bernilai $[-\infty,\infty]$
yang terukur dan terintegralkan.}

\section{*136}{Teorema Kelas Monoton}{22.6.05}{84}{30}
{Aljabar-$\sigma$ yang dibangkitkan oleh suatu keluarga $\Cal{I}$;
aljabar himpunan.}

\medskip

 Lampiran Jilid 1

\chapintrosection{10.10.96}{89}{32}

\section{1A1}{Teori himpunan}{5.11.03}{89}{32}
{Notasi;  himpunan terhitung dan tak terhitung.}

\section{1A2}{Himpunan terbuka dan tertutup dalam $\BbbR^r$}{21.11.03}{92}{33}
{Definisi;  sifat-sifat dasar himpunan terbuka dan tertutup;  ketaksamaan
Cauchy;  bola terbuka.}

\section{1A3}{Limit superior dan limit inferior}{18.12.03}{94}{34}
{$\limsup_{n\to\infty}a_n$, $\liminf_{n\to\infty}a_n$ dalam
$[-\infty,\infty]$.}

\wheader{}{6}{2}{2}{100pt}
%\medskip

Konkordansi \pagereference{97}{}

\medskip

Referensi untuk Jilid 1 \vtmpb{4.10.96}\pagereference{97}{35}

\medskip
     
Indeks Jilid 1
     
\qquad Topik dan hasil utama \pagereference{98}{}
     
\qquad Indeks umum \pagereference{99}{}
     
%102 pages
